\subsection{Joint traffic and execution measurements}\label{sec:cover-study}
The study compares six cover policies on matched synthetic support tickets. Two fixed-weight models, Qwen3:8B and Llama3.1:8B, each run 300 fixtures per condition: 120 training, 60 validation and 120 test, balanced across six actions. The 3,600 episodes have terminal results. Fixtures and cover draws are paired across conditions; request phrasings recur across splits.

\paragraph{Execution and observation}
The models choose checks, business actions, objects, recovery and termination through LangGraph~\cite{langchain2025langgraph}. A controller fetches the ticket and enforces prerequisites without supplying the correct action. Seven local MCP Streamable HTTP destinations use TLS~1.3/HTTP/2, a common \path{/mcp} path and request/response padding quanta of 4,096/16,384 bytes. The observer receives destinations and bidirectional ciphertext byte counts within supplied episode and call windows. Timing and model traffic are outside this view.

A frozen schedule inserts one or two random decoys among early tool calls. The conditions are no cover, proxy discard, read cover, isolated execution, live write, and isolated execution with visible responses. Proxy discard computes on a disposable copy; isolated execution retains a separate clone. Live writes deliberately modify the primary fixture, including other orders or shared quota. The visible control sends clearly marked cover responses from the clone to the planner. Other cover responses are hidden. The live-write arm is an adverse control selected to exercise collateral changes. Proxy discard also computes a disposable response; it is not a pure transport-only dummy.

\paragraph{Matched endpoints and attacks}
\emph{Goal completion} checks the requested action/object, prerequisites and target postcondition. \emph{State integrity} checks the entire business state against the single intended-action reference. We also count completed tasks with state errors directly; the two predicates are not nested. Recognition predicts the requested action including failed executions. Logistic regression (LR) and random forest (RF) train within each condition, with validation selecting the primary attacker. Frozen and destination-position attacks are retained alongside them.

Each endpoint has a prespecified twelve-test Holm family: five cover-minus-baseline contrasts and visible-minus-hidden for each model. Conditional intervals resample paired fixtures by action, holding fitted attacks fixed. Supplement S8 contains every attack and contrast, state replay and a worked evaluation protocol.

\begin{table*}[t]\centering\small
\caption{Executed cover: all held-out conditions, 120 paired fixtures per model and condition. Recognition columns count correct requested-action predictions; chance is 20/120. The learner is selected using validation LR/RF accuracy. The position rule remains visible when it is stronger. Goal, whole-state integrity and completed tasks with state errors are separate endpoints.}
\label{tab:cover-main}
\setlength{\tabcolsep}{4pt}
\begin{tabularx}{\textwidth}{@{}lYrrrrrr@{}}\toprule
Model & Cover policy & Selected & Last dest. & Goal & State & Goal + error & Call bytes $\times$ \\\midrule
Qwen3:8B & No cover & 118 & 116 & 114 & 116 & 0 & 1.00 \\
Qwen3:8B & Proxy discard & 62 & 77 & 114 & 116 & 0 & 1.34 \\
Qwen3:8B & Read cover & 62 & 77 & 114 & 115 & 0 & 1.34 \\
Qwen3:8B & Isolated, hidden & 61 & 77 & 114 & 115 & 0 & 1.34 \\
Qwen3:8B & Live write & 62 & 75 & 102 & 7 & 95 & 1.34 \\
Qwen3:8B & Isolated, visible & 63 & 77 & 113 & 114 & 1 & 1.35 \\
\midrule
Llama3.1:8B & No cover & 113 & 79 & 82 & 69 & 13 & 1.00 \\
Llama3.1:8B & Proxy discard & 79 & 67 & 80 & 68 & 13 & 1.29 \\
Llama3.1:8B & Read cover & 72 & 67 & 84 & 71 & 14 & 1.28 \\
Llama3.1:8B & Isolated, hidden & 82 & 67 & 83 & 70 & 13 & 1.28 \\
Llama3.1:8B & Live write & 80 & 65 & 76 & 3 & 73 & 1.29 \\
Llama3.1:8B & Isolated, visible & 81 & 68 & 83 & 73 & 11 & 1.29 \\
\bottomrule\end{tabularx}
\end{table*}


\paragraph{Completion and collateral effects diverge}
Under live writes, 95 of 102 completed Qwen tasks and 73 of 76 completed Llama tasks have collateral state errors. The corresponding no-cover counts are zero of 114 and thirteen of 82. Whole-state integrity falls from 116/120 to 7/120 and from 69/120 to 3/120. Both paired reductions survive correction; completion falls by 10.0 and 5.0 percentage points, with only Qwen's contrast significant. The adverse writes strongly predispose the direction of state loss. The counts quantify disagreement between the target-task predicate and the broader fixture contract. They neither discover a weakness in state-aware benchmarks nor estimate ordinary cover harms.

\paragraph{Residual recognition and execution cost}
For Qwen, hidden isolated execution completes the same 114/120 tasks as no cover, with state integrity 115/120 versus 116/120 and 1.34 times the tool-call bytes. The highest observed score across the tested attacks falls from 118/120 to 77/120. That latter score comes from the last-destination rule, above the selected learner's 61/120. It exposes the schedule's unprotected tail and gives a more demanding account of residual recognition. At this budget, Qwen's selected attack scores are 62, 62, 61 and 62 for discarded, read, isolated and live-write cover. Live writes show no additional recognition reduction in these point estimates. These observations do not establish equivalence among policies. Across both models, all ten primary cover-versus-baseline recognition contrasts survive correction; substantial recognition remains.

Visible and hidden isolated responses are identical within test pairs. Visibility increases mean prompt tokens by 18.1\% and 18.9\%, while all primary recognition, goal and state contrasts have adjusted $p=1$. The experiment leaves their average effect unresolved. Paired fixtures remain matched despite planner path variation; without repeated runs of each condition, that variation cannot be separated into a stable visibility effect and execution noise. Supplement S8 reports the paths.
